% mazzi: 14
\chapter{Introduzione al futsal}
\label{cap:introfutsal}

\textbf{Modello di prestazione del calcio a 5}: discende dalle \textbf{regole} e
dall'\textbf{evoluzione del gioco}, presupposto della valutazione funzionale, della profilazione di
giocatore e squadra e della programmazione del processo di allenamento.

% ---------------------------------------------------------------------------
\section{Il percorso del modulo}

\begin{enumerate}
  \item \textbf{introduzione al futsal}: la disciplina, il regolamento di gioco e le basi tecniche e
    tattiche;
  \item \textbf{match analysis}, su due piani:
    \begin{itemize}
      \item \textbf{cinematico}: distanze percorse in gara, numero di accelerazioni, profilo del
        movimento;
      \item \textbf{fisiologico}: carico interno durante la partita, rilevato con la frequenza
        cardiaca, più uno studio internazionale con prelievo di lattato;
    \end{itemize}
  \item \textbf{profilo del giocatore di futsal}: quali test di valutazione siano più specifici per
    il calcio a 5 secondo la letteratura scientifica internazionale, con i relativi valori di
    riferimento;
  \item \textbf{valutazione e allenamento del metabolismo anaerobico}, centrato sulla
    \textbf{repeated sprint ability} (RSA), capacità fondamentale nel calcio a 5;
  \item \textbf{valutazione e allenamento del metabolismo aerobico}, centrato sul \textbf{FIET}
    (\textit{Futsal Intermittent Endurance Test}).
\end{enumerate}
Gli ultimi due titoli dicono ``valutazione ed allenamento'' perché il test non si esaurisce in sé:
serve a stabilire i \textbf{parametri dell'allenamento} per la qualità che ha indagato. Si passa
cioè dalla valutazione funzionale al lavoro sul campo.

% ---------------------------------------------------------------------------
\section{Che cos'è il calcio a 5}

\rubrica{Terminologia}
I termini calcio a 5 e futsal si alternano:
\begin{itemize}
  \item \textbf{futsal} è il termine internazionale e deriva dal Brasile, dalla denominazione
    \textit{football de sala};
  \item \textbf{calcio a 5} è il nome usato in Italia fin dai primi tempi, tanto che l'organo
    federale che organizza i campionati si chiama \textbf{Divisione Calcio a 5}.
\end{itemize}

\rubrica{La disciplina}
\begin{itemize}
  \item è uno \textbf{sport indoor}, che si gioca nei palazzetti su una superficie rigida;
  \item presenta aspetti in comune con il calcio: tecnica e tattica individuale, principi di tattica
    collettiva e relativi sviluppi di gioco --- vale per esempio anche la semplice triangolazione
    fra due giocatori in situazione di superiorità numerica;
  \item ha però caratteristiche proprie che rendono specifica la disciplina: è il
    \textbf{regolamento di gioco}, e in particolare alcuni suoi punti, a caratterizzarlo e a
    determinare i comportamenti tecnico-tattici dei giocatori in campo.
\end{itemize}

\rubrica{La storia in Italia}
\begin{enumerate}
  \item \textbf{fine anni Settanta -- inizio anni Ottanta}: nasce a Roma, nei circoli romani, dove
    si giocava sui campi da tennis togliendo la rete;
  \item con il progresso tecnologico dei \textbf{manti sintetici in erba} si passa al ``campetto da
    calcetto'', con dimensioni variabili da una struttura sportiva all'altra;
  \item questo modello non corrispondeva però a ciò che era già un dato di fatto in \textbf{Spagna}
    e in \textbf{Brasile}: uno sport da palazzetto, al chiuso, su superficie rigida --- cemento,
    plastica, linoleum a seconda del livello di qualificazione e delle possibilità economiche dei
    club.
\end{enumerate}

\rubrica{Diffusione}
Sport molto diffuso in Italia sia a livello amatoriale sia a livello organizzato, con un riscontro
in crescita di anno in anno per pubblico nei palazzetti e per visibilità televisiva: l'ultimo
Europeo UEFA è stato trasmesso in Italia da un'emittente satellitare importante, con riscontro di
pubblico molto positivo.

\rubrica{L'organizzazione dei campionati}
\begin{itemize}
  \item \textbf{Serie A}: girone unico, con \textit{play-off} scudetto a fine stagione;
  \item \textbf{Serie A2}: due gironi (centro-nord e centro-sud) fino alla stagione 2017/2018, tre
    dalla 2018/2019;
  \item \textbf{Serie B}: ancora a carattere nazionale, divisa in gironi che comprendono due o tre
    regioni;
  \item \textbf{C1} e \textbf{C2}: campionati regionali;
  \item \textbf{Serie D}: campionato provinciale.
\end{itemize}
A seconda del livello di qualificazione sono richiesti standard diversi: dai campionati nazionali in
su scattano l'obbligo di giocare indoor e i vincoli sulle \textbf{dimensioni del campo}.

\rubrica{Il movimento femminile}
\begin{itemize}
  \item in Italia il calcio a 5 femminile esiste da molti anni, ma ha fatto passi da gigante
    nell'ultimo quinquennio;
  \item il campionato di Serie A è nato nella stagione 2011/2012; il \textbf{girone unico} è però
    recentissimo, in vigore da due stagioni;
  \item è oggi molto seguito, con profili di ottimo livello sia per qualità sia per numero di
    tesserate, e con un bacino in crescita fra le più piccole;
  \item in Spagna lo stesso percorso è avvenuto prima che in Italia: il movimento femminile vi è
    quindi assai più sviluppato.
\end{itemize}

% ---------------------------------------------------------------------------
\section{Il regolamento che rende specifica la disciplina}

\subsection{Spazio, materiali e superfici}

\rubrica{Lo spazio di gioco e il numero dei giocatori}
\begin{itemize}
  \item dimensioni ideali: 40\,m di lunghezza $\times$ 20\,m di larghezza, obbligatorie nelle
    competizioni internazionali;
  \item in Italia non tutte le strutture dispongono di campi di queste dimensioni: se ne trovano di
    lunghezza e larghezza diverse;
  \item \textbf{conseguenza}: visti gli spazi ristretti, giocare su un campo più lungo o più largo
    di 2--5\,m fa differenza sul piano tecnico-tattico, e probabilmente anche su quello
    fisico-atletico --- ma su quest'ultimo punto non esistono studi;
  \item in campo ci sono complessivamente 10 giocatori, cinque per squadra.
\end{itemize}

\rubrica{Il pallone}
Pallone numero 4, a \textbf{rimbalzo controllato}: lasciato cadere a terra rimbalza, ma meno di un
pallone da calcio.

\rubrica{Le superfici di gioco}
\begin{itemize}
  \item all'aperto: sintetico in erba;
  \item indoor: cemento, linoleum, plastica, gomma;
  \item \textbf{superficie ideale}: il parquet.
\end{itemize}

\subsection{Lo svolgimento della gara}

\rubrica{Le rimesse laterali}
Si eseguono con i piedi, ed è stato così esclusivamente per gli ultimi vent'anni. È una regola in
evoluzione: il \textit{board} FIFA si è riunito per rivedere alcune regole del calcio a 5 e renderlo
ancora più spettacolare, e fra le proposte c'è quella di lasciare al giocatore la scelta fra piedi e
mani.

\rubrica{I tiri liberi}
\begin{itemize}
  \item \textbf{che cosa sono}: la possibilità di calciare in porta senza opposizione
    dell'avversario, con palla ferma su un dischetto posto a 10\,m;
  \item \textbf{quando si ottengono}: ogni squadra dispone di 5 falli per tempo; dal sesto in poi
    $\rightarrow$ tiro libero invece della punizione nel punto del fallo;
  \item \textbf{perché}: per calmierare i contatti, che in un campo di 40 $\times$ 20\,m con dieci
    giocatori sarebbero molto frequenti, con rischio di infortuni da trauma;
  \item \textbf{effetti}: la regola ha ridotto molto i falli e ha modificato l'approccio tattico ---
    raggiunto il quinto fallo bisogna limitare i contatti, e cambia il modo di marcare l'avversario.
\end{itemize}

\rubrica{Le sostituzioni volanti}
In qualsiasi momento della partita un giocatore può uscire dal campo, essere sostituito e
\textbf{rientrare} in un secondo momento; al contrario del calcio, dove il giocatore sostituito
chiude lì la propria gara.

\rubrica{La durata della gara}
\begin{itemize}
  \item due tempi da 20 minuti effettivi: il cronometro si ferma ogni volta che la palla non è in
    gioco;
  \item nelle categorie nazionali sono previsti due arbitri in campo, uno per metà campo, più un
    \textbf{cronometrista} a bordo campo incaricato di bloccare il tempo;
  \item \textbf{durata reale}: circa 38--40 minuti il primo tempo e 40--42 minuti il secondo.
\end{itemize}

\rubrica{L'assenza del fuorigioco}
È una delle differenze fondamentali con il calcio.

\rubrica{La ripresa veloce del gioco (4 secondi)}
Qualunque sia il motivo per cui la palla va rimessa in gioco --- fallo laterale, calcio d'angolo ---
il giocatore che deve battere ha 4 secondi, spesso contati ad alta voce dagli arbitri; scaduti
$\rightarrow$ il pallone passa alla squadra avversaria.

\rubrica{I limiti del retropassaggio al portiere}
\begin{itemize}
  \item il portiere può toccare la palla una sola volta nella fase di possesso;
  \item dopo che l'ha giocata, i compagni non possono più ridargli il pallone finché questo non
    viene toccato da un avversario o non entra in possesso della squadra avversaria;
  \item \textbf{conseguenza}: soluzioni tattiche specifiche sia in fase di possesso sia in fase di
    non possesso.
\end{itemize}

\rubrica{Le espulsioni temporanee}
\begin{itemize}
  \item il giocatore espulso esce definitivamente dal terreno di gioco, esattamente come nel calcio:
    ``temporanea'' non si riferisce a lui;
  \item è l'\textbf{inferiorità numerica} a essere temporanea, e dura al massimo 2 minuti: si gioca
    5 contro 4;
  \item gol della squadra in superiorità numerica $\rightarrow$ parità ristabilita subito, anche se
    i due minuti non sono trascorsi.
\end{itemize}

\rubrica{Il time out}
Un \textbf{time out per tempo} a disposizione dell'allenatore: strumento utile sia per correggere le
problematiche tecniche e tattiche, sia per spezzare il ritmo della partita.

% ---------------------------------------------------------------------------
\section{Il futsal come base formativa}

\rubrica{La citazione di Pelé}
``\textit{Futsal requires you to think and play fast. It makes everything easier when you later
switch to football}'': il calcio a 5 richiede di pensare e giocare velocemente, e rende tutto più
facile quando poi si passa al calcio.

\rubrica{Perché è propedeutico al calcio}
\begin{enumerate}
  \item il mondo del calcio sta riconoscendo le forti \textbf{propedeuticità} del futsal nella fase
    iniziale di approccio allo sport, cioè con i bambini;
  \item il motivo principale è la possibilità di semplificare l'insegnamento dei \textbf{gesti
    tecnici fondamentali} --- il passaggio, la conduzione della palla;
  \item il fondamento è \textbf{pedagogico}: per insegnare un movimento si parte dal semplice per
    arrivare al complesso, dal facile per arrivare al difficile;
  \item da qui l'evidenza pratica: è più semplice controllare un pallone che rimbalza poco su una
    superficie piatta e regolare che un pallone che rimbalza molto su una superficie instabile come
    un campo in erba.
\end{enumerate}

\rubrica{Il modello brasiliano}
In Brasile i settori giovanili svolgono allenamenti veri e propri di calcio a 5 --- non allenamenti
di calcio su un campo ristretto --- con istruttori e allenatori di calcio a 5. Non è un caso che i
giocatori brasiliani abbiano concetti di \textbf{tattica individuale} particolarmente ben formati,
primo fra tutti lo \textbf{smarcamento}: è un patrimonio appreso nel calcio a 5.

\rubrica{Il caso del Cruzeiro Esporte Clube}
Programma del settore giovanile di un importante club di Serie A brasiliana, in numero di
allenamenti settimanali:

\begin{table}[htbp]
  \centering
  \small
  \renewcommand{\arraystretch}{1.3}
  \begin{tabularx}{\textwidth}{|l|*{5}{>{\centering\arraybackslash}X|}}
    \hline
    & \textbf{Under 10} & \textbf{Under 11} & \textbf{Under 12} & \textbf{Under 13} &
      \textbf{Under 14} \\
    \hline
    \textbf{Futsal} & 3 & 2 & 1 & 1 & 1 \\
    \hline
    \textbf{Calcio} & 0 & 1 & 3 & 4 & 5 \\
    \hline
  \end{tabularx}
\end{table}

\begin{itemize}
  \item fino all'under 10 --- e a maggior ragione nelle categorie inferiori --- si fanno solo
    allenamenti di calcio a 5, mai di calcio;
  \item dall'under 11 si introduce un allenamento di calcio su campo grande, mantenendo due sedute
    di futsal indoor su parquet;
  \item la progressione arriva a sei sedute settimanali all'under 14, essendo un club
    professionistico;
  \item al termine del percorso il giocatore, consigliato dai tecnici del club, sceglie se
    proseguire nel calcio a 5 o passare al calcio.
\end{itemize}

% ---------------------------------------------------------------------------
\section{Che cosa si vede in campo}

Dalla finale dell'ultimo Europeo, Portogallo--Spagna:

\begin{itemize}
  \item \textbf{il conteggio dei falli}: un segnale acustico avvisa che una squadra ha raggiunto i
    cinque falli; dal successivo in poi l'avversario batte dal dischetto del tiro libero;
  \item \textbf{non esiste il centrocampo}: le squadre o attaccano la porta avversaria o difendono
    la propria. Sono esattamente le situazioni cruciali del calcio, ma qui si ripetono di continuo:
    in ambito giovanile questo produce un numero elevatissimo di esperienze di quel tipo, ed è uno
    stimolo forte all'apprendimento;
  \item \textbf{il vincolo sul portiere in azione}: il portiere può giocare la palla proveniente da
    un fallo laterale, ma da quel momento la sua squadra non può più servirlo finché il pallone non
    viene toccato dagli avversari;
  \item \textbf{i duelli 1 contro 1} si ripetono di continuo: è lo stimolo naturale per
    l'insegnamento di un fondamentale tecnico come il \textbf{dribbling};
  \item \textbf{le situazioni da fermo sono schematizzate}: praticamente tutte le squadre hanno da
    tre a cinque schemi --- o più --- per falli laterali, calci d'angolo e punizioni;
  \item \textbf{il portiere di movimento}: un giocatore che ha tutti i privilegi del portiere ---
    fra cui giocare con le mani nella propria area --- ma che viene schierato per poter attaccare 5
    contro 4 nella metà campo avversaria. Il vantaggio in fase di possesso è indubbio, ma palla
    recuperata dagli avversari $\rightarrow$ porta sguarnita, ed è la prima cosa che cercheranno.
\end{itemize}
